\chapter{Conclusions}
\label{chap:conclusions}

This chapter discusses what the answers to the research questions imply for the development of coor\-di\-na\-tor-based multi-agent systems: how much a probe on the hidden states of the coordinator adds to the routing decision, and how it can be used as a misrouting flag. It then describes the conditions under which a multi-agent system can use such a probe, which concern whether the inference service that runs the coordinator returns its hidden states together with the response, the cost of extracting them, and the labeled data needed to train the probe. Next, the chapter states the limitations of the study, which arise from the synthetic dataset, the structure of the multi-turn conversations, the model pool, and the position from which the probe reads the hidden states. Finally, the chapter outlines directions for future work and closes with the conclusions drawn from the findings of this work.

\section{Implications}
\label{sec:ch9-implications}

The answer to \hyperref[rq:ch8-rq1]{\textbf{RQ1}} in Section~\ref{subsec:ch8-rq1-answer} shows that a probe on the hidden states of a coordinator selects the correct sub-agent more reliably than the coordinator itself. The probe is a linear classifier on the hidden state at the last prompt token, which the coordinator already computes, so a multi-agent system obtains a more accurate routing decision without an additional model call. The gain grows as the correct sub-agent becomes more difficult to determine from the wording of the request, which is where the coordinator misroutes most often. The gain is also smaller when the request is part of a multi-turn conversation.

The answer to \hyperref[rq:ch8-rq2]{\textbf{RQ2}} in Section~\ref{subsec:ch8-rq2-answer} shows that the same gain is not obtained by the three text-only classifiers of this work, which are logistic regressions trained on a bag-of-words, a TF-IDF and an embedding representation of the request text and need no access to the hidden states. As the correct sub-agent becomes more difficult to determine from the wording of the request, the accuracy of each text-only classifier falls more than the accuracy of the probe. On these requests, the hidden states of the coordinator therefore carry routing information that the representations of the request text do not provide, and this information is what justifies reading the hidden states. As the difficulty of the requests increases, the probe is the classifier that retains the most accuracy to evaluate the routing decision of the coordinator.

The answer to \hyperref[rq:ch8-rq3]{\textbf{RQ3}} in Section~\ref{subsec:ch8-rq3-answer} shows that the optimal layer for the probe falls at a consistent relative depth across the models of this work. For a new coordinator model, it is therefore likely to fall at a similar depth, and the layer sweep aimed at finding it can be restricted around that depth.

Because the probe is more accurate than the coordinator, it can be used as an evaluation flag on the routing decision of the coordinator. The answer to \hyperref[rq:ch8-rq4]{\textbf{RQ4}} in Section~\ref{subsec:ch8-rq4-answer} shows that the disagreement flag, raised when the probe and the coordinator select different sub-agents, marks the requests that the coordinator is likely to misroute, with a precision above the base error rate of the coordinator in every run. When the flag is raised, the execution of the request can be stopped (Section~\ref{sec:early-exit-feedback-loop}), or the request can be kept for later review by a human. On the single-turn dataset, the flag is raised on between 56.0\% and 79.7\% of the misroutes of the coordinator, depending on the model, and between 68.1\% and 82.6\% of the flagged requests are misroutes of the coordinator. On the multi-turn dataset, the flag is raised on between 54.0\% and 77.8\% of the misroutes, and between 59.8\% and 80.5\% of the flagged requests are misroutes of the coordinator.

Because the probe is more accurate than the coordinator, the request could also be handed off to the sub-agent that the probe predicts, instead of the sub-agent that the coordinator selects. The probe would then replace the routing decision of the coordinator, instead of serving as an evaluation flag on it. With this replacement, a smaller coordinator can give a more accurate routing decision than a larger one: on both datasets, the probes on Qwen3 4B Instruct 2507 and on Gemma 4 E2B IT are more accurate than the Llama 3.3 70B Instruct and the Qwen2.5 72B Instruct coordinators (Tables~\ref{tab:ch7-qwen3-4b}, \ref{tab:ch7-gemma4-e2b-it}, \ref{tab:ch7-llama33-70b} and~\ref{tab:ch7-qwen25-72b}). Ratnakar et al. reach the same conclusion with LatentGate, where a linear probe on the hidden states of Qwen2.5 1.5B Instruct routes queries more accurately than prompt-based routers built on GPT-4.1, GPT-4.1-nano and Gemini 2.5 Flash \cite{ratnakarLatentGateLowLatencySemantic2026}. However, this replacement may not be possible when the coordinator has responsibilities other than routing, such as conducting the conversation with the user, producing the final response from the response of the selected sub-agent (Section~\ref{sec:coordinator}), or applying guardrails (Section~\ref{sec:security-evaluation}).


\section{Conditions of Use}
\label{sec:ch9-conditions}

The implications just described assume that the hidden states of the coordinator are available and that a probe has been trained on them. A multi-agent system meets this assumption only under certain conditions of use, which concern whether the inference server returns the hidden states of the coordinator, the memory needed to extract them, and the labeled data needed to train the probe.

LLM-based agents commonly interact with large language models through an \emph{inference server}, a process that loads a model and runs it each time it receives a request. The model runs on the GPUs of the machine that hosts the server, so the agent does not need to load the model weights or run the forward pass itself. The agent invokes the model through the API endpoints that the server exposes: it sends the prompt in a request and receives the generated text in the response. These endpoints can belong to a frontier-model API that a provider runs on its own infrastructure, such as the OpenAI API or the Anthropic API. They can also belong to a self-hosted inference server, deployed on self-managed infrastructure. Common self-hosted inference servers include \texttt{llama.cpp}, \texttt{vLLM}, \texttt{SGLang}, and \texttt{MLX}.

\subsection{Hidden States in the Server Response}
\label{subsec:ch9-flag-access}

Frontier-model APIs do not provide access to the hidden states. Exposing them may also be a security risk for the provider, because they allow a third party to reconstruct parts of a proprietary model without access to its weights. Carlini et al. showed that the log-probabilities and the logit bias exposed by production APIs were sufficient to recover the final projection layer and the hidden dimension of OpenAI models \cite{carliniStealingPartProduction2024}. The hidden states expose more of the model than these outputs. Jiao et al. used \emph{knowledge distillation} to train a smaller model that reproduces the hidden states of a larger model, so that the smaller model could be used in place of the larger one \cite{jiaoTinyBERTDistillingBERT2020}. A third party that receives the hidden states of a proprietary model could apply the same method to obtain a smaller model that replaces the proprietary one.

Self-hosted inference servers have no standard method to provide access to the hidden states. Each server that supports this access implements it in a different way. The documentation of \texttt{llama.cpp} and of \texttt{MLX} lists no option to return the hidden states \cite{LlamacppToolsServer, MlxlmMlx_lmSERVERmd}. The \texttt{vLLM} server can extract the hidden states of the prompt tokens at selected layers, but it does not return them to the caller: it writes them to a file on disk and returns only the path of the file in the response \cite{teamExtractingHiddenStates2026}. Among the four servers, only \texttt{SGLang} returns the hidden states in the response, through a server argument that enables this feature \cite{ServerArguments}.

\subsection{Inference Cost}
\label{subsec:ch9-flag-cost}

Where the hidden states are accessible, extracting them increases the memory that inference requires. As Section~\ref{subsec:hidden-states} describes, the hidden state $h^{(\ell)}_i$ of every layer $\ell$ and every token $i$ is a vector of $d_{\mathrm{model}}$ values. The extracted hidden states are held in VRAM or RAM, and the memory they occupy grows linearly with the number of layers and tokens they cover.

Each value of a hidden state vector occupies $b$ bytes, with $b = 2$ for 16-bit and $b = 4$ for 32-bit floating-point values. The hidden states of one token therefore occupy $b \times d_{\mathrm{model}}$ bytes at one layer and $b \times L \times d_{\mathrm{model}}$ bytes at all $L$ layers. Table~\ref{tab:ch9-hidden-state-size} reports both sizes for the models of Table~\ref{tab:model-properties}.

\begin{table}[htbp]
    \centering
    \small
    \begin{tabular}{lrrrrrr}
        \toprule
        & & & \multicolumn{2}{c}{One layer (KiB)} & \multicolumn{2}{c}{All $L$ layers (KiB)} \\
        \cmidrule(lr){4-5} \cmidrule(lr){6-7}
        Model & $d_{\mathrm{model}}$ & $L$ & 16-bit & 32-bit & 16-bit & 32-bit \\
        \midrule
        Gemma 4 E2B IT & 1536 & 35 & 3.0 & 6.0 & 105.0 & 210.0 \\
        Qwen3 4B Instruct 2507 & 2560 & 36 & 5.0 & 10.0 & 180.0 & 360.0 \\
        Qwen3 8B & 4096 & 36 & 8.0 & 16.0 & 288.0 & 576.0 \\
        Gemma 4 12B IT & 3840 & 48 & 7.5 & 15.0 & 360.0 & 720.0 \\
        GPT-OSS 20B & 2880 & 24 & 5.6 & 11.3 & 135.0 & 270.0 \\
        Gemma 4 26B A4B IT & 2816 & 30 & 5.5 & 11.0 & 165.0 & 330.0 \\
        Llama 3.3 70B Instruct & 8192 & 80 & 16.0 & 32.0 & 1280.0 & 2560.0 \\
        Qwen2.5 72B Instruct & 8192 & 80 & 16.0 & 32.0 & 1280.0 & 2560.0 \\
        \bottomrule
    \end{tabular}
    \caption[Hidden-state memory per token, at one layer and at all $L$ layers]{Memory occupied by the hidden states of one token, at one layer and at all $L$ layers, with 16-bit and 32-bit floating-point values (1~KiB = 1024 bytes).}
    \label{tab:ch9-hidden-state-size}
\end{table}

Requesting the hidden states together with the response from an inference server may therefore increase the memory usage of the server and slow down inference. It may also considerably increase the size of the payload that the server returns for each request.

\subsection{Cold-Start Problem}
\label{subsec:ch9-flag-cold-start}

A probe is subject to a \emph{cold-start problem}, because it needs labeled data to be trained. Every probe in this work is trained on samples whose correct sub-agent is set by the generation plan of Chapter~\ref{chap:dataset-construction}. A real production system typically has no labeled data at the time of its first release. Labeled data accumulates only after the system is deployed, either from human annotation or from usage history annotated after the routing decision. Once enough labeled data has accumulated, a probe can be trained and used as a misrouting flag. However, the probe would need to be retrained every time the routing task changes, for example when a sub-agent is added, removed, or redefined. It would also need to be retrained every time the coordinator model changes, because a different model may have a different optimal layer, and the accuracy gained by probing that layer may also differ. A change of the coordinator prompt may also affect the accuracy of the probe, because the hidden state at the last prompt token is computed from the entire prompt. The effect of a prompt change on the accuracy of the probe is not measured in this work, and neither is the number of labeled samples a probe needs before it can be used as a misrouting flag. Section~\ref{sec:ch9-future-work} lists both questions as future work.


\section{Limitations of the Study}
\label{sec:ch9-limitations-work}

The implications that Section~\ref{sec:ch9-implications} draws from the results of Chapter~\ref{chap:results-analysis} are subject to the limitations stated in this section, which concern the synthetic dataset, the structure of the multi-turn conversations, the scale and composition of the model pool, and the position from which the probe reads the hidden states.

\subsection{Synthetic Dataset}
\label{subsec:ch9-limitations-dataset}

Every result reported in this work is measured on the single-turn and multi-turn datasets described in Chapter~\ref{chap:dataset-construction}. Both datasets are generated synthetically by prompting a large language model, and each sample is built to meet the requirements stated in Section~\ref{sec:dataset-requirements}. The content and the distribution of the requests are therefore determined by the design of the generation pipeline. Every accuracy in this work is measured on this distribution. The requests that a deployed multi-agent system receives from its users may follow a different distribution. The accuracy gain of the probe over the coordinator in such a system may therefore be larger or smaller than the one measured in Chapter~\ref{chap:results-analysis}.

No human annotation is performed on the datasets. The label of each sample comes from the generation plan, and it is validated only by the automated audit stage. As Section~\ref{subsec:pipeline-audit-repair} describes, this stage applies deterministic structural rules, the Jaccard index on \texttt{adversarial} pairs, and the verdict of an LLM oracle. A manual revision is applied only to the samples that the oracle check flags, so every other sample keeps the label assigned by the plan without human verification. Some of these samples may therefore be mislabeled, which would affect every accuracy reported in Chapter~\ref{chap:results-analysis}.

\subsection{Structure of the Multi-Turn Conversations}
\label{subsec:ch9-limitations-multiturn}

The generation pipeline builds every multi-turn sample with the same structure. As Section~\ref{subsec:multiturn-derivation} describes, each conversation contains five messages, which alternate between the developer and the coordinator. In the \texttt{adversarial} tier, the two members of a pair share the first four messages, so that they differ only in the final message.

A conversation between a developer and a coordinator may span tens or hundreds of turns. Its topic may change between turns, and messages unrelated to the routing decision may be interleaved with the messages that carry the information needed for routing.

On the multi-turn dataset, the probe loses more accuracy than the coordinator in all seven models, as Section~\ref{subsec:ch8-multi-turn} reports. Because every conversation has the same structure and the same number of turns, this drop is observed only on conversations of this form, and the dataset does not allow an analysis of how the probe behaves on conversations of different lengths and structures. It is therefore unknown whether the probe remains more accurate than the coordinator on such conversations, and whether the disagreement flag can still mark the misroutes of the coordinator.

\subsection{Scale and Composition of the Model Pool}
\label{subsec:ch9-limitations-scale}

As Section~\ref{sec:model-choice} describes, every model in the pool is an open-weight model with at most 72 billion parameters. The model pool is therefore limited to this parameter count, and the results of this work apply only to models up to this size.

The model pool is also limited in composition, because its models are not distributed evenly across families and sizes. Of the seven models analyzed in Chapter~\ref{chap:results-analysis}, three belong to the Gemma 4 family and two to the Qwen3 family, and all five have at most 26 billion parameters. Llama 3.3 70B and Qwen2.5 72B are the only two models above that size, and each is the only model of its family in the pool. The results are consistent across the four families, but they are derived mostly from models of the Gemma 4 and Qwen3 families, so they do not show whether the accuracy gain of the probe over the coordinator depends on the model family.

\subsection{Probe Input Position}
\label{subsec:ch9-limitations-extraction}

As Section~\ref{subsec:probe-position-requirements} states, every probe in this work reads the hidden state at the last prompt token. The only exception is GPT-OSS 20B, whose hidden state at the last reasoning token is also probed for comparison in Section~\ref{sec:ch8-gpt-oss}. The last prompt token is chosen because of masked attention: as Section~\ref{subsec:last-token} describes, it is the only position whose hidden state is computed from the entire prompt.

The position from which the hidden states are read may affect the accuracy of a probe. For GPT-OSS 20B, neither the last prompt token nor the last reasoning token gives a probe that improves on the coordinator, so the position must be chosen according to the model architecture. For example, Healy et al. reported that detection accuracy varies with the method used to extract hidden states from a query \cite{healyInternalRepresentationsIndicators2026}. LatentGate is one example of a different choice, because it pools the hidden states over every token of the query instead of reading a single position \cite{ratnakarLatentGateLowLatencySemantic2026}. Since every probe in this work is trained only on the hidden state at the last prompt token, the results do not allow a comparison between different probing positions.


\section{Future Work}
\label{sec:ch9-future-work}

\paragraph{Validation against production data.} On the datasets of this work, the probe is more accurate than the coordinator in seven of the eight models, and its disagreement with the coordinator marks coordinator errors, as Section~\ref{subsec:ch8-rq4-answer} states. Every sample of these datasets is generated from a fixed taxonomy of five sub-agents and follows the patterns defined by the generation plan. Requests collected from a coordinator-based multi-agent system in use may not follow a predefined set of patterns, and they may be incomplete, ambiguous, or unrelated to any sub-agent. A labeled dataset of such requests would measure the accuracy of the probe under this variability, and this accuracy would determine how useful the probe is as a misrouting flag in a multi-agent system in use.

A multi-agent system in use may already have a set of production requests annotated with the correct sub-agent, which is used to evaluate the routing decision of its coordinator. Future work could validate the probe on such a system with the following procedure:

\begin{enumerate}
    \item \textbf{Training.} The hidden states of the coordinator are collected by running each annotated request through the prefill step (Section~\ref{subsec:prefill-decode}) of the model in use, with no response token generated. A probe is trained on these hidden states with $k$-fold cross-validation, as for the probes of Section~\ref{subsec:ch7-probe}.
    \item \textbf{Deployment.} The probe is deployed next to the coordinator, on an inference server that returns the hidden states (Section~\ref{subsec:ch9-flag-access}), and it raises the disagreement flag on the incoming requests.
    \item \textbf{Collection period.} The flagged requests are collected for a defined period of time.
    \item \textbf{Annotation.} At the end of the collection period, human annotators examine every flagged request and determine whether the coordinator misrouted it or the flag is a false positive.
\end{enumerate}

The fraction of the flagged requests that are misrouted gives the precision of the flag in real usage. This precision would show how effective a probe trained on evaluation data is as a misrouting flag on the requests of a multi-agent system in use.

\paragraph{Adaptation to a changing sub-agent set.} The sub-agent set of a multi-agent system in use may change over time, for example when a new sub-agent is added or when the responsibilities of an existing one are redefined. Section~\ref{subsec:ch9-flag-cold-start} notes that a probe needs retraining after each of these changes. The labeled data collected before the change may no longer match the new routing task, so the probe could not be used as a misrouting flag until enough new labeled data has accumulated. A change of the sub-agent set also modifies the prompt of the coordinator, because every sub-agent is declared in the prompt as a tool (Section~\ref{subsec:ch7-coordinator-prompt}). The hidden state at the last prompt token is computed from the entire prompt, so it may change even for requests whose correct sub-agent stays the same. The hidden states may also change when the prompt of the coordinator is modified and the sub-agent set stays the same.

Future work could study how the accuracy of an existing probe varies after each type of change: the addition, the removal, and the redefinition of a sub-agent. This study would show which changes leave the accuracy of the probe unaffected and which ones require new labeled data. Future work could also study whether an existing probe can be adapted to the new sub-agent set by reusing what it has already learned, for example by updating it with only the samples labeled after the change. If the adapted probe needs fewer labeled samples than a new probe to reach the same accuracy, fewer requests would need to be labeled after each change before the probe can be used again as a misrouting flag.

\paragraph{Detailed analysis of multi-turn behavior.} On the multi-turn dataset, the multi-turn drop of the probe is larger than the multi-turn drop of the coordinator in all seven models, as Section~\ref{subsec:ch8-multi-turn} reports. Section~\ref{subsec:ch9-limitations-multiturn} states that every conversation in this dataset has the same structure and the same number of messages. A conversation in a multi-agent system in use may span many more turns, change topic, and include turns unrelated to the routing decision. The hidden state at the last prompt token is computed from the entire conversation, so the information needed for routing may be spread across more turns or placed far from the last prompt token. The accuracy of the probe on such conversations may therefore differ from the one measured on the multi-turn dataset of this work.

Future work could build a synthetic multi-turn dataset in which the number of turns and the structure of the conversation vary as design factors. The structure could follow a defined taxonomy of patterns, such as topic drift, turns unrelated to the routing decision, and different positions of the turn that carries the routing information. Such a dataset would allow a study of how the accuracy of the probe varies across the patterns, and across shorter and longer conversations. The same dataset could also be used to test whether a probe that reads the hidden states at several positions, such as the last token of each turn, is more accurate on long conversations, or on conversations with topic drift, than a probe that reads only the last prompt token. The results would show whether the probe remains usable as a misrouting flag on conversations longer and less regular than the ones of this work.


\section{Final Remarks}
\label{sec:ch9-conclusion}

On requests whose correct sub-agent is difficult to determine from their wording, a linear probe on the hidden state at the last prompt token selects the correct sub-agent more often than the coordinator in seven of the eight models, and more often than the text-only classifiers. Its margin over the coordinator is largest on the requests that the coordinator misroutes most often. The hidden states of a coordinator therefore carry routing information that its stated decision does not reflect and that the wording of the request does not provide. The exception is GPT-OSS 20B, which reasons before it selects a sub-agent. Its last prompt token comes before the point where the routing decision may be formed, so the probing position must be chosen according to the model architecture.

On multi-turn conversations, the probe loses more accuracy than the coordinator in all seven models, so its margin is smaller. The routing information at the last prompt token may therefore be weaker when the request is spread over several turns.

The disagreement between the probe and the coordinator marks the requests that the coordinator is likely to misroute, with a precision above the base error rate of the coordinator in all fourteen runs, although it raises more false positives on multi-turn conversations. When the hidden states of the coordinator are accessible, a flagged request can be stopped or held for human review before it affects the rest of the run.

These results indicate that a linear probe on the hidden states of a coordinator can provide an evaluation signal on its routing decision. Such a probe may support failure attribution and improve the quality of coordinator-based multi-agent systems, because it identifies likely routing failures before they propagate through the rest of the multi-agent trace.

